% =============================================================================
% Chapter 7: Debugging Lessons and Fixes
% =============================================================================
\chapter{Debugging Lessons and Fixes}
\label{ch:debugging}

This chapter documents the major issues encountered during testbench development and how each
was resolved. These lessons are essential for anyone extending the testbench, porting it to a
different simulator, or building testbenches for similar APB peripherals.

\section{Issue Summary Table}

\begin{table}[H]
\centering
\caption{Summary of Issues Encountered and Resolved}
\label{tab:issues}
\small
\begin{tabularx}{\textwidth}{@{}clXX@{}}
\toprule
\textbf{\#} & \textbf{Issue} & \textbf{Symptom} & \textbf{Resolution} \\
\midrule
1 & SDI signal type & RX always read 0x00 & Changed \texttt{logic} to \texttt{wire} for SDI signals \\
2 & Loopback vs. force & SDI not reaching DUT RX & Used \texttt{force dut.spi\_sdi1} for diagnostic \\
3 & SPI clock lag & Slave data unstable at sample point & Used negedge-clocked counter BFM \\
4 & APB read timing & RXFIFO returned wrong word & Sampled PRDATA at posedge before \texttt{\#1} delay \\
5 & RXFIFO visibility & Could not verify RX data in FIFO & Added hierarchical probes into \texttt{dut.u\_rxfifo} \\
6 & EOT synchronization & EOT missed in testbench & Implemented sticky \texttt{eot\_flag} latch \\
7 & Spurious interrupt & Potential GitHub-reported bug & Confirmed FSM resets to INACTIVE (safe) \\
\bottomrule
\end{tabularx}
\end{table}

\section{Issue 1: SDI Signal Driving in Xcelium}

\subsection{Problem}

The first testbench version declared all SPI data-in signals as \texttt{logic} and used
\texttt{assign} statements to create loopback connections:

\begin{lstlisting}
logic spi_sdi0, spi_sdi1, spi_sdi2, spi_sdi3;
assign spi_sdi0 = spi_sdo0;  // Loopback
assign spi_sdi1 = spi_sdo1;  // Loopback
\end{lstlisting}

In Xcelium, this \texttt{assign} on a \texttt{logic} signal \emph{silently fails}. The signals
remain stuck at their initial value (0), causing all RX reads to return \texttt{0x00000000}.

\subsection{Root Cause}

In SystemVerilog, \texttt{logic} is a single-driver type. An \texttt{assign} statement constitutes
a continuous driver, but Xcelium treats this as a conflict with the \texttt{logic} declaration and
silently suppresses the driver without issuing a warning.

\subsection{Solution}

Changed all SDI signals to \texttt{wire} type with inline continuous assignments:
\begin{lstlisting}
wire spi_sdi0 = spi_sdo0;
wire spi_sdi1 = slave_drive_en ?
     slave_tx_data[5'd31 - slave_bit_cnt] : spi_sdo1;
wire spi_sdi2 = spi_sdo2;
wire spi_sdi3 = spi_sdo3;
\end{lstlisting}

\subsection{Lesson Learned}

Always use \texttt{wire} for signals driven by \texttt{assign} statements in Xcelium. This
applies to any testbench signal that uses continuous assignment, including loopback connections,
BFM outputs, and multiplexed signals.

\section{Issue 2: Testbench-Side Loopback vs. DUT-Internal Forcing}

\subsection{Problem}

Even after fixing the wire type issue, TC4 initially failed because the testbench-side
\texttt{spi\_sdi1} assignment was not guaranteed to be visible at the DUT's internal sampling
point due to delta-cycle ordering.

\subsection{Solution}

For diagnostic purposes, TC4 uses \texttt{force dut.spi\_sdi1 = 1'b1} to bypass all testbench
wiring and inject a known value directly at the DUT port. This isolates the RX datapath test
from any BFM or wiring issues.

After TC4 confirmed the RX path worked, the hardware-style slave BFM was developed for
production-quality data injection (used in TC6, TC8, TC10).

\subsection{Lesson Learned}

When debugging data-integrity failures, use hierarchical \texttt{force} as a diagnostic tool
to isolate the DUT from the testbench. Once the DUT is proven correct, incrementally add
testbench complexity (BFM, protocol models) and re-verify.

\section{Issue 3: SPI Clock Lag and Slave BFM Timing}

\subsection{Problem}

The DUT's \texttt{spi\_clk} output is a registered version of the internal clock. The internal
\texttt{spi\_rise} event (used by the RX shift register) precedes the external \texttt{spi\_clk}
rising edge by one system clock cycle. A task-based slave BFM clocked on
\texttt{@(posedge spi\_clk)} would present new data \emph{after} the DUT has already sampled,
resulting in the DUT reading stale data.

\begin{verbatim}
System clock:    _/‾\_/‾\_/‾\_/‾\_/‾\_
spi_rise (int):  ____/‾\________________  (internal, 1 cycle early)
spi_clk (ext):   ________/‾\____________  (registered, 1 cycle late)
Slave data chg:  ____________/‾\________  (TOO LATE for DUT sample!)
\end{verbatim}

\subsection{Solution}

Replaced the task-based BFM with a counter-based \texttt{always\_ff} block clocked on
\texttt{negedge spi\_clk}. The bit counter increments on falling edges, and the data output
is combinationally derived from the counter. This ensures data is stable at least half an SPI
clock period before the next rising edge.

\subsection{Lesson Learned}

When building slave BFMs for SPI masters, always account for the clock registration delay.
Use falling-edge clocking for data transitions, and prefer combinational data indexing over
sequential data shifting to avoid race conditions.

\section{Issue 4: APB Read Timing for RXFIFO}

\subsection{Problem}

The RXFIFO's \texttt{data\_o} output is combinational: \texttt{data\_o = buffer[pointer\_out]}.
The \texttt{pointer\_out} register increments via non-blocking assignment when
\texttt{ready\_i \& valid\_o} is true. In the APB read task, if \texttt{PRDATA} was sampled
\emph{after} the \texttt{\#1} delay (postponed region), the pointer had already advanced, and
\texttt{data\_o} showed the \emph{next} FIFO entry.

\subsection{Solution}

In the \texttt{apb\_read} task, \texttt{PRDATA} is sampled at the \texttt{posedge HCLK}
\emph{before} the \texttt{\#1} delay:
\begin{lstlisting}
@(posedge HCLK);     // Access phase posedge
rdata = PRDATA;      // Sample in Active region, BEFORE pointer update
#1;                  // Now pointer_out may advance
PSEL = 1'b0;         // Deassert
\end{lstlisting}

\subsection{Lesson Learned}

For any APB peripheral with FIFO read ports, sample \texttt{PRDATA} at the clock edge (Active
region), not after a delay. This is because the FIFO's read pointer advances on the same edge
via non-blocking assignment, and the combinational data output immediately reflects the new
pointer value.

\section{Issue 5: RXFIFO Visibility and Hierarchical Probing}

\subsection{Problem}

Early RX failures showed \texttt{0x00000000} from \texttt{pop\_rx()}, but it was unclear
whether the problem was in the RX shift register, the FIFO push logic, the FIFO read logic,
or the APB read path.

\subsection{Solution}

TC4 adds extensive hierarchical probes to observe internal state at three key time points
(pre-trigger, during transfer, post-EOT):

\begin{lstlisting}
$display("PRE-TRIGGER: ctrl state=%0d, rx_CS=%0d, fifo_elem=%0d",
         dut.u_spictrl.state,
         dut.u_spictrl.u_rxreg.rx_CS,
         dut.u_rxfifo.elements);

$display("POST-EOT: rx data_int=0x%08X, fifo_elem=%0d",
         dut.u_spictrl.u_rxreg.data_int,
         dut.u_rxfifo.elements);

$display("POST-EOT: fifo buffer[0]=0x%08X, ptr_out=%0d",
         dut.u_rxfifo.buffer[0],
         dut.u_rxfifo.pointer_out);
\end{lstlisting}

This revealed that \texttt{data\_int} was correctly 0xFF, \texttt{elements} was 1, and
\texttt{buffer[0]} contained 0xFF---proving the RX path was correct and the issue was in the
APB read timing.

\subsection{Lesson Learned}

When debugging data-integrity issues, add hierarchical probes at every stage of the datapath.
Check: (1) the input signal at the DUT port, (2) the shift register contents, (3) the FIFO
element count, (4) the FIFO buffer contents directly, (5) the FIFO pointer values, and
(6) the APB read data.

\section{Issue 6: EOT Synchronization}

\subsection{Problem}

The \texttt{events\_o[1]} EOT pulse is only one system clock cycle wide. A testbench that
polls EOT in a loop (\texttt{while(!events\_o[1]) @(posedge HCLK)}) can miss the pulse if
the poll and the pulse align to the same edge.

\subsection{Solution}

The testbench uses a sticky \texttt{eot\_flag} latch (see Chapter~\ref{ch:tb_arch}) that
captures the EOT pulse and holds it until explicitly cleared. The \texttt{wait\_eot} task
polls this flag, not the raw \texttt{events\_o[1]} signal.

\subsection{Lesson Learned}

For any single-cycle event signal, use a sticky latch in the testbench rather than
polling the raw signal. This eliminates timing-dependent test failures.

\section{Issue 7: Spurious Interrupt After Enable}

\subsection{Problem}

A GitHub issue reported that enabling interrupts after reset could cause a spurious interrupt
pulse, because the interrupt condition (TX FIFO empty $\leq$ threshold 0) is immediately true.

\subsection{Investigation}

TC14 was written to reproduce this bug. It monitors \texttt{events\_o[0]} for 20 cycles before
and 5 cycles after writing \texttt{int\_en = 1}.

\subsection{Finding}

The bug was \emph{not present} in the current RTL. The interrupt FSMs initialize to
\texttt{INT\_xx\_INACTIVE} on reset. In this state, the FSM does not check the threshold
condition---it waits for either an INTSTA read (non-counter mode) or a counter target match
(counter mode) to transition to \texttt{INT\_xx\_ACTIVE}. Only in the ACTIVE state does
the threshold condition trigger an interrupt.

\subsection{Lesson Learned}

Always verify reported bugs against the actual RTL version in use. Write a targeted testcase
that reproduces the expected failure mode, and document the result whether or not the bug
is confirmed.

\section{Best Practices to Avoid Regressions}

Based on the debugging experience, the following practices should be followed for future work:

\begin{enumerate}[leftmargin=2em]
  \item \textbf{Use \texttt{wire} for all driven signals} in Xcelium testbenches.
  \item \textbf{Sample APB read data at posedge}, before any \texttt{\#} delay.
  \item \textbf{Use sticky latches} for single-cycle event signals.
  \item \textbf{Add hierarchical probes} when debugging data-integrity failures.
  \item \textbf{Use falling-edge BFMs} for SPI slave models.
  \item \textbf{Always wait for EOT} before issuing software reset or reconfiguring.
  \item \textbf{Write regression testcases} for every bug found, even if already fixed.
  \item \textbf{Keep the global watchdog} to catch FSM lockups.
\end{enumerate}
